\section*{Bab \ref{ch:g4:volcanoes-earthquakes-first-look} --- Gunung Api dan Gempa Bumi: Pengenalan Pertama}

\begin{solution}{exo:g4:volcanoes-earthquakes-first-look:1}
(a) \omterm{def:g4:volcanoes-earthquakes-first-look:eruption}{Lava}. (b) \omterm{def:g4:volcanoes-earthquakes-first-look:eruption}{Abu vulkanik}. (c) Gas.
\end{solution}

\begin{solution}{exo:g4:volcanoes-earthquakes-first-look:2}
\omterm{def:g4:volcanoes-earthquakes-first-look:earthquake}{Gempa bumi} adalah guncangan tanah yang tiba-tiba, yang disebabkan oleh
\omterm{def:g3:stones-pebbles-sand:rock}{batuan} yang patah dan bergeser jauh di bawah tanah. \omterm{def:g4:volcanoes-earthquakes-first-look:earthquake}{Gempa bumi} berlangsung
beberapa detik (satu atau dua menit untuk yang paling besar), bukan
berjam-jam.
\end{solution}

\begin{solution}{exo:g4:volcanoes-earthquakes-first-look:3}
Merunduk, berlindung, berpegangan.
\end{solution}

\begin{solution}{exo:g4:volcanoes-earthquakes-first-look:4}
Di sekeliling \omterm{def:g2:water-on-earth:ocean}{Samudra} Pasifik.
\end{solution}

\begin{solution}{exo:g4:volcanoes-earthquakes-first-look:5}
(a) Benar: kebanyakan gempa hanya dirasakan oleh alat ukur. (b) Salah: \omterm{def:g4:volcanoes-earthquakes-first-look:eruption}{lava}
mendingin dan mengeras menjadi \omterm{def:g3:stones-pebbles-sand:rock}{batuan} baru. (c) Salah: berlari ketika tanah
berguncang itu berbahaya (kamu bisa jatuh, dan kaca serta benda-benda
berjatuhan dari gedung); merunduklah, berlindunglah, dan berpeganganlah.
\end{solution}

\begin{solution}{exo:g4:volcanoes-earthquakes-first-look:6}
\omterm{def:g4:volcanoes-earthquakes-first-look:eruption}{Lava} berasal dari jauh di bawah \omterm{def:g4:volcanoes-earthquakes-first-look:volcano}{gunung api}, tempat \omterm{def:g3:stones-pebbles-sand:rock}{batuan} meleleh. \omterm{def:g4:volcanoes-earthquakes-first-look:eruption}{Lava} naik
melalui bagian tengah \omterm{def:g4:volcanoes-earthquakes-first-look:volcano}{gunung api} dan keluar di puncaknya.
\end{solution}

\begin{solution}{exo:g4:volcanoes-earthquakes-first-look:7}
(a) Chili (Peru, Ekuador, atau Meksiko juga benar): \omterm{def:g4:volcanoes-earthquakes-first-look:volcano}{gunung-gunung apinya}
berjajar di sepanjang negeri itu, dekat \omterm{def:g1:landscapes-around-us:water}{pesisir} Pasifik. (b) Hawaii, di
tengah \omterm{def:g2:water-on-earth:ocean}{Samudra} Pasifik, jauh dari cincin yang diarsir.
\end{solution}

\begin{solution}{exo:g4:volcanoes-earthquakes-first-look:8}
$141 \div 12$ sedikit kurang dari 12 (karena $12 \times 12 = 144$): kira-kira
12 \omterm{def:g4:volcanoes-earthquakes-first-look:earthquake}{gempa bumi} besar setiap bulan.
\end{solution}

\begin{solution}{exo:g4:volcanoes-earthquakes-first-look:9}
Dari 20, seperempatnya adalah 5, jadi tiga perempatnya adalah $3 \times 5 = 15$:
kira-kira 15 di dalam cincin, dan kira-kira $20 - 15 = 5$ di luarnya.
\end{solution}

\begin{solution}{exo:g4:volcanoes-earthquakes-first-look:10}
Seharusnya ia merunduk dengan bertumpu pada tangan dan lututnya, melindungi
kepala dan lehernya di bawah mejanya, lalu berpegangan pada meja itu sampai
guncangan berhenti. Ambang pintu tidak melindunginya dari benda yang jatuh
atau terlempar, dan berlari ketika tanah berguncang bisa membuatnya jatuh.
\end{solution}

\begin{solution}{exo:g4:volcanoes-earthquakes-first-look:11}
Di luar cincin: $1214 - 837 = 377$ \omterm{def:g4:volcanoes-earthquakes-first-look:volcano}{gunung api}. Sepertiga dari 1\,200 adalah
400, jadi dua pertiga dari 1\,200 adalah 800. Cincin itu memuat 837, sedikit
lebih dari 800 dari kira-kira 1\,200: ``kira-kira 2 dari 3'' adalah gambaran
yang tepat.
\end{solution}
